\subsection{群胚被群作用的商}

设 G 为一个传递群胚，K 为一个群，并设 \(\theta\colon\mathrm{K}\to\mathrm{Aut}(\mathrm{G})^{\circ}\) 为从 K 到群胚 G 的自同构群之相反群中的群同态。我们将说群 K 右作用于 G。若 \(k\in\mathrm{K}\) ，我们有时以 \(k^{*}x\) （相应地 \(k^{*}f\) ）表示 G 的顶点 \(x\) （相应地箭 \(f\) ）在群胚自同构 \(\theta(k)\) 下的像。

设 \(|\mathrm{K}|\) 为这样的群胚：其顶点集为 \(\mathrm{K}\) ，连接两个顶点的箭集当这两个顶点不同时为空集，否则缩为单个元素。设 \(\mathrm{H}\) 为积群胚 \(\mathrm{G} \times |\mathrm{K}|\) ；\(\mathrm{H}\) 的一个顶点是偶 \((a,k)\) ，其中 \(a\) 是 \(\mathrm{G}\) 的一个顶点，\(k\) 是 \(\mathrm{K}\) 的一个元素；若 \(f\) 是 \(\mathrm{G}\) 中连接顶点 \(a\) 与顶点 \(b\) 的一条箭，我们以 \((f,k)\) 表示 \(\mathrm{H}\) 中连接 \((a,k)\) 与 \((b,k)\) 的唯一的箭。设 \(\varphi: \mathrm{H} \to \mathrm{G}\) 为由第一个投影给出的群胚态射，并设 \(\psi: \mathrm{H} \to \mathrm{G}\) 为满足 \(\operatorname{Som}(\psi)((a,k)) = k^*a\) 与 \(\operatorname{Fl}(\psi)((f,k)) = k^*f\) （若 \(k \in \mathrm{K}\) 、\(a \in \operatorname{Som}(\mathrm{G})\) 且 \(f \in \operatorname{Fl}(\mathrm{G})\) ）的群胚态射。

以 G/K 表示余均衡子 \(\operatorname{Coeg}(\varphi,\psi)\) ，并以 \(\gamma\colon G\to G/K\) 表示典范群胚态射。

设 o 为 G 的一个顶点。对 \(k \in K\) ，选取 G 中的一条箭 \(c_{k}\) （连接 \(k^{*}o\) 与 o 的）；由于 G 是传递的，这样的箭存在。对 \(k \in K\) ，以 \(\text{Fix}(k)\) 表示 G 的子群胚，其顶点（相应地箭）是被 k 固定的 \(\text{Som}(G)\) （相应地 \(\text{Fl}(G)\) ）的元素；我们选取一个顶点集 \(A_{k}\) （\(\text{Fix}(k)\) 的，与此群胚的所有轨道都相交）。对 \(k \in K\) 及 \(a \in A_{k}\) ，我们还选取 G 中连接顶点 a 与顶点 o 的一条箭 \(f_{(a,k)}\) 。

命题 6. —— 唯一的群同态

\[
\lambda \colon \mathrm{G} _ {o} * \mathrm{F} (\mathrm{K}) \to (\mathrm{G} / \mathrm{K}) _ {\gamma (o)}
\]

使得 \(\lambda(f) = \gamma_o(f)\) （对 \(f \in \mathrm{G}_o\) ）以及 \(\lambda([k]) = \gamma(c_k)\) （对 \(k \in \mathrm{K}\) ）者是满射的。其核是 \(\mathrm{G}_o * \mathrm{F}(\mathrm{K})\) 中包含下列元素的最小正规子群：

\[
r _ {2} (k, f) = [ k ] ^ {- 1} f [ k ] (c _ {k} ^ {- 1} k ^ {*} (f) ^ {- 1} c _ {k})\tag{\((R_{2})\)}
\]

\[
\text { for } k \in \mathrm{K} \text { and } f \in \mathrm{G} _ {o};\tag{\((\mathrm{R}_{3}^{\prime})\)}
\]

\[
r _ {3} ^ {\prime} (k, a) = [ k ] (c _ {k} ^ {- 1} k ^ {*} (f _ {(a, k)}) ^ {- 1} f _ {(a, k)}))
\]

\[
\text { for } k \in \mathrm{K} - \{e \} \text { and } a \in \mathrm{A} _{k}\tag{\((\mathrm{R}_{3}^{\prime\prime})\)}
\]

\[
r _ {3} ^ {\prime \prime} (k, h) = [ k h ] ^ {- 1} [ k ] [ h ] (c _ {h} ^ {- 1} h ^ {*} (c _ {k} ^ {- 1}) c _ {k h})
\]

\[
\text { for } k,h \in \mathrm{K}.
\]

由于 G 是传递的，把元素 \(k \in K\) 对应到 \((o, k)\) 的轨道的映射是 K 到 H 的轨道集上的一个双射。我们于是把 \(\operatorname{Orb}(\mathbf{G})\) 与 \(\{o\}\) 等同，把 \(\operatorname{Orb}(\mathbf{H})\) 与 K 等同。框架 \(\Gamma_{0}\) （偶 \((\varphi, \psi)\) 的）于是等同于具有单个顶点 o 且箭集为 K 的箭图。设 T 为 \(\Gamma_{0}\) 的、顶点集为 \(\{o\}\) 而箭集为空的子箭图；其相伴图是图 \(\widetilde{\Gamma}_{0}\) 的唯一的极大树。

族 \((o,(o,k)_{k\in \mathrm{K}},(e_o)_{k\in \mathrm{K}},(c_k)_{k\in \mathrm{K}},\mathrm{T},o)\) 是偶 \((\varphi ,\psi)\) 的一个基本配备。

映射 \(f \mapsto (f, k)\) 定义了迷向群 \(G_o\) 到迷向群 \(H_{(o,k)}\) 上的一个同构；在此同构下，同态 \(\varphi_k\) 与 \(\psi_k\) （从 \(H_{(o,k)}\) 到 \(G_o\) 中的、由 II, 第 208 页的命题 5 定义的）由下式给出

\[
\varphi_ {k} (f, k) = f \quad \text { and } \quad \psi_ {k} (f, k) = c _ {k} ^ {- 1} k ^ {*} (f) c _ {k},\tag{3}
\]

（对 \(k\in \mathrm{K}\) 及 \(f\in \mathrm{G}_o\) ）。

箭图 \(H \times_{G} H\) 是 \(H \times H\) 中箭图态射 \(\psi \circ pr_{1}\) 与 \(\varphi \circ pr_{2}\) 的均衡子。其顶点是偶 \(((a, k), (b, h))\) ，其中 a 与 b 是 G 的顶点，k 与 h 是 K 中满足 \(b = k^{*}a\) 的元素；其箭是偶 \(((f, k), (g, h))\) ，其中 f 与 g 是 G 的箭，k 与 h 是 K 中满足 \(g = k^{*}f\) 的元素；箭 \(((f,k),(g,h))\) 的起点等于 \(((o(f),k),(o(g),h))\) ；其靶等于 \(((t(f),k),(t(g),h))\) 。

我们随后定义一个箭图态射 \(\mu\) （从 \(H \times_{\mathrm{G}} H\) 到 H 中的），令 \(\mu((a,k),(b,h)) = (a,kh)\) 与 \(\mu((f,k),(g,h)) = (f,kh)\) 。我们有 \(\varphi \circ \mu = \varphi \circ \mathrm{pr}_1\) 与 \(\psi \circ \mu = \psi \circ \mathrm{pr}_2\) 。

设 \(x\) 与 \(y\) 为 H 中满足 \(\varphi(x) = \varphi(y)\) 的顶点。于是存在 G 的一个顶点 \(a\) 以及 K 的元素 \(k\) 与 \(h\) ，使得 \(x = (a, k)\) 且 \(y = (a, h)\) 。令 \(z = (h^* a, h^{-1} k)\) ；我们有 \(\mu(y, z) = x\)。这表明偶 \((\varphi, \psi)\) 满足 II, 第 202 页的命题 4 的假设。

H 的顶点 \((a,k)\) 属于群胚 \(\operatorname{Ker}(\varphi,\psi)\) （\(\varphi\) 与 \(\psi\) 的均衡子）当且仅当 \(k^{*}a=a\) ，即 k 属于群 K 中顶点 a 的稳定子。设 \(A_{1}\) 为偶 \((a,k)\) （对 \(k\in K\) 与 \(a\in A_{k}\)）所成的集；它与 H 的子群胚 \(\operatorname{Ker}(\varphi,\psi)\) 的所有轨道相交。设 \(Z_{1}\) 为 \(\Omega(\widetilde{\Gamma})\) 中由形如 (((a,k),1)) （对 \((a,k)\in A_{1}\)）的序列组成的部分。对 \(\mathbf{z}=((a,k),1)\in\mathbf{Z}_{1}\) ，\((f_{(a,k)},k)\) 是 H 的一条箭，它连接 \((a,k)\) 与 \((o,k)\) 。令 \(f(\mathbf{z})=((f_{(a,k)},k))\)。

顶点所成的集 \(\mathrm{A}_2\) （\(\mathrm{H} \times_{\mathrm{G}} \mathrm{H}\) 的，由形如 \(((o,k), (k^*o,h))\) （对 \((k,h) \in \mathrm{K}^2\) ）的顶点组成）与 \(\mathrm{H} \times_{\mathrm{G}} \mathrm{H}\) 的每个轨道相交。还可注意到 \(\mu((o,k), (k^*o,h)) = (o,kh)\) 。H 中的箭 \((e_o,k), (c_k,h), (e_o,kh)\) 分别连接 \((o,k), (k^*o,h), (o,kh)\) 到 \((o,k), (o,h)\) 与 \((o,kh)\) 。以 \(\mathrm{Z}_2\) 表示形如 \((((o,k),1), ((k^*o,h), 1), ((o,kh), -1))\) 的序列在 \(\Omega(\widetilde{\Gamma})\) 中所成的集；对这样一个元素 \(\mathbf{z}\) （\(\mathrm{Z}_2\) 的），令 \(f(\mathbf{z}) = ((e_o,k), (c_k,h), (e_o,kh))\) 。

由 II, 第 202 页的命题 4，集 \(Z = Z_1 \cup Z_2\) 与族 \((f(\mathbf{z}))_{z \in \mathbb{Z}}\) 构成一个补配备。

设 \(k \in \mathrm{K}\) 且 \(a \in \mathrm{A}_k\) 。元素 \(\mathrm{C}((a, k), 1)\) （由 II, 第 208 页的公式 (2) 定义的、群 \(\mathrm{G}_o * \mathrm{F}(\mathrm{K})\) 中的）等于

\[
c _ {k} ^ {- 1} \psi ((f _ {(a, k)}, k)) ^ {- 1} \varphi (f _ {(a, k)}) e _ {o} [ k ] = (c _ {k} ^ {- 1} k ^ {*} (f _ {(a, k)}) ^ {- 1} f _ {(a, k)}) [ k ].\tag{4}
\]

设 \(k,h \in \mathbf{K}\)。我们验证，元素

\[
\mathrm{C} (((o, k), 1), ((k ^ {*} o, h), 1), ((o, k h), - 1))
\]

（由 II, 第 208 页的公式 (2) 定义的、群 \(G_{o} * F(K)\) 中的）等于

\[
[ k ] [ h ] (c _ {h} ^ {- 1} h ^ {*} (c _ {k} ^ {- 1}) c _ {k h}) [ k h ] ^ {- 1}.\tag{5}
\]

元素 \(r_3'(e, a)\) （由关系式 \((\mathrm{R}_3')\) 对 \(k = e\) 给出的）都等于 \([e]c_e^{-1}\) ，后者是群 \(\mathrm{G}_o * \mathrm{F}(\mathrm{K})\) 中把关系式 \((\mathrm{R}_{3}^{\prime\prime})\) 应用于 \(k = h = e\) 而得到的元素。鉴于关系式 (3)、(4) 与 (5)，命题于是由 II, 第 208 页, 命题 5 得出。

推论 1. —— 假设群 K 由 G 的顶点的固定子生成。那么群同态 \(\gamma_{o}\colon\mathrm{G}_{o}\to(\mathrm{G}/\mathrm{K})_{\gamma(o)}\) 是满射的。此外，若群胚 G 是单传递的，则群胚 G/K 亦然。

关系式 \((\mathrm{R}_{3}^{\prime\prime})\) 蕴含 \(\lambda([e]) = \lambda(c_{e}) = \gamma_{o}(c_{e})\) 。关系式 \((\mathrm{R}_{3}^{\prime\prime})\) 进而蕴含：那些 \(k \in K\) （满足 \(\lambda([k])\) 属于 \(\gamma_{o}\) 的像的）组成 K 的一个子群。最后，关系式 \((\mathrm{R}_{3}^{\prime})\) 表明，对每个固定子非空的元素 \(k \in K\) ，\(\lambda([k])\) 属于 \(\gamma_{o}\) 的像。推论由此得出。

注记 1. —— 我们可以给出群 \((\mathrm{G}/\mathrm{K})_{\gamma(o)}\) 的另一个有时更方便的描述。为此，令 \(M = K \times G_{o}\) ，并用下列公式在 M 上定义一个合成法则

\[
(k, a) \cdot (h, b) = (k h, c _ {k h} ^ {- 1} h ^ {*} (c _ {k} a) c _ {h} b),
\]

（对 \(k\) 、\(h \in \mathrm{K}\) 及 \(a\) 、\(b \in \mathrm{G}_o\) ）。我们验证这个合成法则是结合的，\((e, c_e^{-1})\) 是一个恒等元，且元素 \((k^{-1}, c_{k^{-1}}^{-1}(k^{-1})^*(c_k a)^{-1} c_e)\) 是 \((k, a)\) 的逆。因此它赋予 M 一个群结构。此外，映射 \(\lambda'\colon \mathrm{M} \to (\mathrm{G} / \mathrm{K})_{\gamma(o)}\) （由 \((k, a) \mapsto \lambda([k]a)\) 定义的）是一个群同态。设 \(\alpha'\) 为从 \(\mathrm{G}_o * \mathrm{F}(\mathrm{K})\) 到 M 的唯一的群同态，满足 \(\alpha'(f) = (e, f)\) （若 \(f \in \mathrm{G}_o\) ）及 \(\alpha'([k]) = (k, e_o)\) （若 \(k \in \mathrm{K}\) ）；我们有 \(\lambda' \circ \alpha' = \lambda\) 。

关系式 \((\mathrm{R}_{2})\) 与 \((\mathrm{R}_{3}^{\prime\prime})\) 表明，\((\mathrm{G}/\mathrm{K})_{\gamma(o)}\) 的每个元素都是同态 \(\lambda\) 下、\(\mathrm{G}_{o}*\mathrm{F}(\mathrm{K})\) 中形如 \([k]f\) 的元素的像，其中 \(f\in G_{o}\) 且 \(k\in K\) 。因此，同态 \(\lambda'\) 是满射的。我们同样验证：\(\alpha'\) 下、\(\mathrm{G}_{o}*\mathrm{F}(\mathrm{K})\) 中形如 \((\mathrm{R}_{2})\) 或 \((\mathrm{R}_{3}^{\prime\prime})\) 的元素的像为零。于是，同态 \(\lambda'\) 的核是 M 的、包含 \(\alpha'\) 下 \(\mathrm{G}_{o}*\mathrm{F}(\mathrm{K})\) 中形如 \((\mathrm{R}_{3}^{\prime})\) 的元素的像的最小正规子群。

推论 2. —— 假设群 K 自由地作用于 Som(G)。那么存在唯一的群同态 \(\pi\colon (\mathrm{G} / \mathrm{K})_{\gamma (o)}\to \mathrm{K}\) ，其核包含 \(\gamma_{o}\) 的像，且满足 \(\pi (\lambda ([k])) = k\) （对所有 \(k\in \mathbf{K}\) ）。此外，\(\mathrm{G}_o\xrightarrow{\gamma_o} (\mathrm{G} / \mathrm{K})_{\gamma (o)}\xrightarrow{\pi}\mathrm{K}\) 是 K 被 \(\mathrm{G}_o\) 的一个扩张。

若这样一个群同态 \(\pi\) 存在，则群同态 \(\pi \circ \lambda\) 必然等于唯一的群同态 \(p\) （从 \(\mathrm{G}_{o} * \mathrm{F}(\mathrm{K})\) 到 \(\mathrm{K}\) 中的），满足 \(p(f) = e\) （对 \(f \in \mathrm{G}_{o}\) ）与 \(p([k]) = k\) 。立即可验证，\(\mathrm{G}_{o} * \mathrm{F}(\mathrm{K})\) 的、由公式 \((\mathrm{R}_{2})\) 与 \((\mathrm{R}_{3}^{\prime \prime})\) 定义的元素属于 \(p\) 的核。根据假设，不存在 \((\mathrm{R}_{3}^{\prime})\) 型的元素。于是，态射 \(\lambda\) 的核包含 \(p\) 的核。因此，存在唯一的群同态 \(\pi: (\mathrm{G} / \mathrm{K})_{\gamma(o)} \to \mathrm{K}\) ，使得 \(\pi \circ \lambda = p\) 。

显然同态 \(\pi\) 是满射的。为证明同态 \(\gamma_{o}\) 是单射的且其像恰为 \(\pi\) 的核，我们注意到同态 \(\lambda'\colon M \to (\mathrm{G}/\mathrm{K})_{\gamma(o)}\) （II, 第 213 页, 注记 1）是一个同构，因为我们已假设 K 自由地作用于 \(\operatorname{Som}(\mathrm{G})\)。复合同态 \((\lambda')^{-1} \circ \gamma_{o}\) （从 \(\mathrm{G}_{o}\) 到 M 中的）由 \(f \mapsto (e, f)\) 给出，而同态 \(\pi \circ \lambda'\colon M \to K\) 把 \((k, f)\) 映为 \(k\)。推论由此得出。

推论 3. —— 假设群胚 G 是单传递的。设 K₀ 为由 G 的顶点的稳定子生成的 K 的子群。从 K 到 \((\mathrm{G}/\mathrm{K})_{\gamma(o)}\) 的、把 \(k\) 映为 \(\gamma(c_k)\) 的映射是一个满射群同态，其核为 K₀。

若一个元素 \(k \in \mathrm{K}\) 固定顶点 \(a\) （\(\mathrm{G}\) 的），则元素 \(g^{-1}kg\) 固定顶点 \(g^{*}a\) ；这蕴含 \(\mathrm{K}_{0}\) 是 \(\mathrm{K}\) 的一个正规子群。

由命题 5，唯一的同态 \(\lambda\colon\mathrm{F}(\mathrm{K})\to(\mathrm{G}/\mathrm{K})_{\gamma(o)}\) （满足 \(\lambda([k])=\gamma(c_k)\) 的）是满射的，且其核是 F(K) 中包含元素 \([k]\) （其中 \(k\) 是固定 G 的一个顶点的、K 的元素）以及元素 \([kh]^{-1}[k][h]\) （其中 \((k,h)\in\mathrm{K}^2\) ）的最小正规子群。特别地，映射 \(\lambda'\colon\mathrm{K}\to(\mathrm{G}/\mathrm{K})_{\gamma(o)}\) （由 \(\lambda'(k)=\gamma(c_k)=\lambda([k])\) 对 \(k\in\mathrm{K}\) 定义的）是一个群同态。我们有 \(\lambda=\lambda'\circ p\) ，其中 \(p\colon\mathrm{F}(\mathrm{K})\to\mathrm{K}\) 表示典范满射群同态。因此，同态 \(\lambda'\) 是满射的，且其核是 K 中包含元素 \(k\) （对固定 G 的一个顶点的 \(k\) ）的最小正规子群，即 \(\mathrm{K}_0\)。命题得证。
% semantic-fragments: legacy-input-seam
\relax{}
